\color{unidarkgreen}{\subsection{Логика отключения}}
\color{black}

\begin{enumerate}

%%%%%%%%%%%%%%%%%%%%%%%%%%%%%%%%%%%%%%%%%%%%%%%%%% ОБщие сведения %%%%%%%%%%%%%%%%%%%%%%%%%%%%%%%%%%%%%%%%%%%%%%%%%%
\item \textbf{Общие сведения}

\begin{enumerate}

\item 
Логика отключения ошиновки состоит из логики отключения и логики запрета АПВ.

\item 
Ввод/вывод функции осуществляется либо посредством программного ключа \textbf{«ЛО СШ»}, либо с использованием внешнего сигнала \textit{«Вывод ЛО СШ»}.

\item 
Уставки логики отключения ошиновки представлены в таблице \ref{app_set:losh}.

\end{enumerate}


%%%%%%%%%%%%%%%%%%%%%%%%%%%%%%%%%%%%%%%%%%%%%%%%%% Логика отключения %%%%%%%%%%%%%%%%%%%%%%%%%%%%%%%%%%%%%%%%%%%%%%%%%%
\item \textbf{Логика отключения}

\begin{enumerate}

\item
Логика отключения предназначена для формирования сигналов отключения ошиновки, а также сигналов отключения и запрета АПВ присоединений. Отключение ошиновки может осуществляться от ДЗО, ЧТО или от УРОВ присоединений.

\item
Для фиксации срабатывания защиты на время отключения всех выключателей, время выполнения цикла АПВ, а также включения всех присоединений участвующих в этом цикле применяется сигнал \textit{«Фикс. откл. СШ»}, который продлевается на выдержку времени \textbf{«Тфикс. откл. СШ»}. Указанный сигнал используется в логике очувствления и запрета АПВ.

\item
Функциональная схема логики отключения ошиновки приведена на рисунке \ref{pic_lo:Log_otkl_ssh}.

\medskip
\begin{figure}[H]
\centering
\includegraphics[width=0.9\textwidth,height=0.9\textheight,keepaspectratio]{\fbpath/05.02 ДЗО/661.9901 Логика отключения/img/Log_otkl_ssh.pdf}
\captionof{figure}{Функциональная схема логики отключения ошиновки}\label{pic_lo:Log_otkl_ssh}
\end{figure}


\end{enumerate}

%%%%%%%%%%%%%%%%%%%%%%%%%%%%%%%%%%%%%%%%%%%%%%%%%% Логика запрета АПВ %%%%%%%%%%%%%%%%%%%%%%%%%%%%%%%%%%%%%%%%%%%%%%%%%%
\item \textbf{Логика запрета АПВ}

\begin{enumerate}

\item
Логика запрета АПВ предназначена для формирования запрета АПВ всех присоединений СШ, чтобы минимизировать вероятность повторного возникновения аварийной ситуации. 

\item
Сигналы общего запрета АПВ формируется в следующих условиях:
\begin{itemize}
\item при неуспешном АПВ;
\item при неполнофазном или полнофазном отказе выключателя питающего присоединения;
\item при срабатывании УРОВ;
\item при неуспешном опробовании;
\item при внешнем отключении;
\item при оперативном запрете АПВ.
\end{itemize}

\item
Неуспешное АПВ выявляется при повторном срабатывании защиты после первого отключения. Для фиксации первого отключения используется сигнал \textit{«Фикс. откл. СШ»}, который формируется в логике
отключения. Для того чтобы сигнал неуспешного АПВ формировался только при повторном срабатывании защиты, используется дополнительная выдержка времени, задаваемая уставкой \textbf{«Т АПВ 1пр. СШ»}. Данная выдержка времени должна охватывать время отключения всех выключателей, но при этом должна быть отстроена от времени включения первого присоединения в цикле АПВ.

\item
Неполнофазный или полнофазный отказ выключателя питающего присоединения выявляется по факту наличия напряжения после сформированного сигнала отключения. Отключение СШ фиксируется при наличии
сигнала \textit{«Фикс. откл. СШ»} и отсутствии сигнала \textit{«Отключение СШ»}. Выдержка времени 25 мс необходима для отстройки от разновременности отключения фаз выключателя.

\item
Запрет АПВ от внешних защит осуществляется от сигналов срабатывания внешнего УРОВ присоединений.

\item
Запрет АПВ при отключении от УРОВ осуществляется при включении программного ключа \textbf{«ЗАПВ от УРОВ СШ»}.

\item
Запрет АПВ при неуспешном опробовании осуществляется при включении программного ключа \textbf{«ЗАПВ при опроб. СШ»}.

\item
Ввод оперативного запрета АПВ при отключении осуществляется внешним сигналом \textit{«Опер. запрет АПВ»}. Оперативный запрет применяется в следующих случаях:
\begin{itemize}
\item при выполнении оперативных переключений с участием разъединителей, когда из-за возможных ошибок может сработать ДЗО и выполнение АПВ недопустимо;
\item после проведения ремонтных работ на ошиновке, когда существует вероятность включения на КЗ.
\end{itemize}

\item
Функциональная схема логики запрета АПВ приведена на рисунке \ref{pic_lo:Log_ZAPV}.

\medskip
\begin{figure}[H]
\centering
\includegraphics[width=0.9\textwidth,height=0.9\textheight,keepaspectratio]{\fbpath/05.02 ДЗО/661.9901 Логика отключения/img/Log_ZAPV.pdf}
\captionof{figure}{Функциональная схема логики запрета АПВ}\label{pic_lo:Log_ZAPV}
\end{figure}

\end{enumerate}


\end{enumerate}